Many federal grants for local public investment reach a place only if a local government applies for them, so places facing similar hazards can end up with very different amounts of investment. This article examines how county size relates to participation in the Hazard Mitigation Grant Program (HMGP), the largest U.S. hazard mitigation grant, after tornado and severe-storm disaster declarations. We link 12,081 county--declaration observations (2,767 counties, 736 declarations, 1999--2021) to HMGP subapplications, including withdrawn and denied requests, and to funded projects. Comparing counties designated in the same declaration, which share the incident, the state administration, and the grant allocation, the smallest third of counties are 8.8 percentage points less likely to obtain any HMGP project, against a mean of 27.1 percent among larger counties, and about half as likely to obtain a tornado safe room or wind retrofit. The shortfall arises at the application stage: once small counties apply, their requests are approved and withdrawn at rates similar to those of larger counties. Fewer eligible local governments, such as municipalities, school districts, and special districts, account for about two-fifths of the gap, but county governments in small counties are themselves less likely to apply. Per-resident fiscal measures do not account for the gap, a 2018 federal provision reimbursing management costs did not detectably narrow it, and the gap is smaller in states whose agencies filed a larger share of applications. The estimates describe associations with county size rather than causal effects of administrative capacity. They identify the application stage as a point at which geographic disparities in access to intergovernmental public investment arise.
